\section{Cruscotto di Valutazione}


%%%%%%%%%%%%%%%%%%%%%%%%%%%%%%%%%%%
%%%%%%%%%%TUTTO DA RIFARE%%%%%%%%%%
%%%%%%%%%%%%%%%%%%%%%%%%%%%%%%%%%%%

\subsection{MQC\_01: Earned Value e MQC\_02: Planned Value}
\begin{figure}[H]
	\centering 
	\includegraphics[width=0.8\textwidth]{PianoQualifica/PB/EarnedPlanned}
	\caption{MQC\_01: Earned Value e MQC\_02: Planned Value}
\end{figure}

Dal grafico emerge la capacità del gruppo nel completare tutte le attività previste in ciascuno sprint$^G$, portando alla coincidenza fra le metriche \vr{Earned Value$^G$} e \vr{Planned Value$^G$}.\\
Si osserva inoltre un progressivo aumento del lavoro pianificato da svolgere, soprattutto negli ultimi sprint in cui i componenti del gruppo non avevano altri impegni universitari.

\subsection{MQC\_03: Actual Cost e MQC\_07: Estimate to Complete}
\begin{figure}[H]
	\centering 
	\includegraphics[width=0.8\textwidth]{PianoQualifica/PB/ActualEstimate}
	\caption{MQC\_03: Actual Cost e MQC\_07: Estimate to Complete}
\end{figure}
Dal grafico si osserva che il gruppo ha gestito bene i costi: i costi effettivi sono sempre rimasti al di sotto della metrica \vr{Estimate at Completion}$^G$ e del budget massimo di 11.440€; solo nell'ultimo sprint il costo effettivo ha superato la stima finale, pur restando comunque entro il budget.
Nella prima fase del progetto$^G$ si può vedere come la differenza fra i preventivi e i consuntivi sia più evidente, tuttavia nel corso del tempo la capacità di pianificazione$^G$ del gruppo è migliorata, riducendo progressivamente questo scarto. Questa evoluzione è documentata nel \href{https://nullpointersgroup.github.io/Documentazione/output/PB/Documenti%20Esterni/Piano_di_Progetto.pdf}{Piano di Progetto$^G$}.

\subsection{MQC\_04: Cost Performance Index e MQC\_05: Schedule Performance Index}
\begin{figure}[H]
	\centering 
	\includegraphics[width=0.8\textwidth]{PianoQualifica/PB/CostSchedulePerformance}
	\caption{MQC\_04: Cost Performance Index e MQC\_05: Schedule Performance Index}
\end{figure}
Dall'analisi si evince che il Cost Performance Index$^G$ (CPI) si mantiene stabilmente al di sopra del valore ottimale, mentre lo Schedule Performance Index$^G$ (SPI) coincide perfettamente. Questo risultato testimonia l'efficacia$^G$ del gruppo nel controllare i costi, mantenendoli sistematicamente al di sotto del budget preventivato, senza mai rischiare di superarlo alla conclusione del progetto.

\subsection{MQC\_06: Estimate at Completion}
\begin{figure}[H]
	\centering 
	\includegraphics[width=0.8\textwidth]{PianoQualifica/PB/EstimateAtCompletion}
	\caption{MQC\_06: Estimate at Completion}
\end{figure}
Il costo stimato di fine progetto$^G$ si attesta sempre sotto la soglia del budget definito in sede di Candidatura.
Il picco del costo stimato è stato raggiunto allo Sprint$^G$ 11, periodo in cui il gruppo era interamente impegnato nella codifica dell'MVP e della stesura dei documenti.

\subsection{MQC\_08: Time Estimate at Completion}
\begin{figure}[H]
	\centering 
	\includegraphics[width=0.8\textwidth]{PianoQualifica/PB/TimeEstimate}
	\caption{MQC\_08: Time Estimate at Completion}
\end{figure}
Poiché lo Schedule Performance Index$^G$ si è sempre mantenuto al valore ottimale di 1, anche la stima del tempo necessario per la conclusione del progetto$^G$ è rimasta invariata, confermando i 176 giorni previsti tra la data di aggiudicazione dell'appalto e la scadenza prevista del 30 aprile.

\subsection{MQC\_09: Completeness Issue}
\begin{figure}[H]
	\centering 
	\includegraphics[width=0.8\textwidth]{PianoQualifica/PB/Completeness}
	\caption{MQC\_09: Completeness Issue}
\end{figure}
Il gruppo nel corso del progetto$^G$ è sempre riuscito a completare nel periodo di sprint$^G$ le attività delineate, evidenziando la capacità di assegnare e distribuire correttamente le attività fra i membri del gruppo.

\subsection{MQC\_10: Indice di Gulpease}
\subsubsection{Indice di Gulpease Analisi dei Requisiti}
\begin{figure}[H]
	\centering 
	\includegraphics[width=0.8\textwidth]{PianoQualifica/PB/GulpeaseAnalisi}
	\caption{Indice di Gulpease Analisi dei Requisiti}
\end{figure}
L'indice di Gulpease dell'Analisi dei Requisiti$^G$ non ha mai superato la soglia del valore ottimo. Questo è probabilmente dovuto al fatto che il documento contiene molte figure, tabelle ed elenchi, elementi che lo script di calcolo esclude in quanto non pertinenti al compito dell'indice stesso.\\
Di conseguenza, il contenuto puramente testuale è molto limitato seppur descrittivo e dettagliato, determinando l'impossibilità strutturale di raggiungere il valore ottimale della metrica.

\subsubsection{Indice di Gulpease Piano di Qualifica}
\begin{figure}[H]
	\centering 
	\includegraphics[width=0.8\textwidth]{PianoQualifica/PB/GulpeaseQualifica}
	\caption{Indice di Gulpease Piano di Qualifica}
\end{figure}
Per le stesse motivazioni esposte in precedenza, anche l'indice di Gulpease del Piano di Qualifica$^G$ non è riuscito a superare la soglia ottimale definita dal gruppo.

\subsubsection{Indice di Gulpease Piano di Progetto}
\begin{figure}[H]
	\centering 
	\includegraphics[width=0.8\textwidth]{PianoQualifica/PB/GulpeaseProgetto}
	\caption{Indice di Gulpease Piano di Progetto}
\end{figure}
L'indice di Gulpease del Piano di Progetto si attesta sopra la soglia dell'ottimalità dal periodo della PB, segno della capacità del gruppo di riuscire a definire in maniera più semplice e immediata il documento.

\subsubsection{Indice di Gulpease Norme di Progetto}
\begin{figure}[H]
	\centering 
	\includegraphics[width=0.8\textwidth]{PianoQualifica/PB/GulpeaseNorme}
	\caption{Indice di Gulpease Norme di Progetto}
\end{figure}
Le \href{https://nullpointersgroup.github.io/Documentazione/output/PB/Documenti%20Interni/Norme_di_Progetto.pdf}{Norme di Progetto$^G$} si sono stabilizzate tra lo sprint 3 e lo sprint 7, periodo in cui il gruppo ha completato la definizione delle principali procedure operative.\\
Successivamente le modifiche sono state minime, come documentato nel \vr{Registro delle Modifiche}.\\
Tuttavia, a seguito della revisione RTB, è stato necessario rivedere il documento per sistemare alcune parti segnalate: questa riscrittura ha fatto aumentare l'indice di 8 punti, per poi rimanere stabile in tutti gli sprint successivi.

\subsubsection{Indice di Gulpease Glossario}
\begin{figure}[H]
	\centering 
	\includegraphics[width=0.8\textwidth]{PianoQualifica/PB/GulpeaseGlossario}
	\caption{Indice di Gulpease Glossario}
\end{figure}
Il Glossario$^G$ contiene per definizione termini specialistici e tecnici, influendo strutturalmente sul suo indice di Gulpease, mantenendolo su valori più bassi.

\subsubsection{Indice di Gulpease Specifica Tecnica}
\begin{figure}[H]
	\centering 
	\includegraphics[width=0.8\textwidth]{PianoQualifica/PB/GulpeaseSpecifica}
	\caption{Indice di Gulpease Specifica Tecnica}
\end{figure}
La Specifica Tecnica fa un uso esteso di diagrammi e figure, indispensabili per rappresentare le relazioni tra i moduli. \\
La parte descrittiva, seppur non predominante, è comunque presente per le necessarie spiegazioni e rispetta gli standard qualitativi del gruppo.

\subsubsection{Indice di Gulpease Manuale Utente}
\begin{figure}[H]
	\centering 
	\includegraphics[width=0.8\textwidth]{PianoQualifica/PB/GulpeaseManuale}
	\caption{Indice di Gulpease Manuale Utente}
\end{figure}
Il Manuale Utente è prevalentemente visivo: le numerose immagini guidano l'utente passo passo, riducendo la necessità di lunghi testi. \\
La parte testuale, seppur limitata, rispetta gli standard qualitativi del gruppo.

\subsection{MQC\_11: Test Success Rate}
\begin{figure}[H]
	\centering 
	\includegraphics[width=0.8\textwidth]{PianoQualifica/PB/TestSuccess}
	\caption{MQC\_11: Test Success Rate}
\end{figure}
All'inizio di ogni sprint, tutti i test scritti dal gruppo superavano la fase senza errori bloccanti, dimostrando la capacità di scrivere test corretti e di correggere prontamente quelli falliti. \\
Tuttavia, questa metrica è piuttosto limitante, poiché non evidenzia con precisione i momenti di maggiore criticità dovuti al refactoring del codice.

\subsection{MQC\_12: Quality Metrics Satisfied}
\begin{figure}[H]
	\centering 
	\includegraphics[width=0.8\textwidth]{PianoQualifica/PB/QualityMetrics}
	\caption{MQC\_12: Quality Metrics Satisfied}
\end{figure}
Nel grafico sono considerate soddisfatte le metriche almeno accettabili. \\
Esso mostra la stabilità del gruppo nel raggiungimento delle metriche.\\
A differenza della RTB, è stata rispettata pure l'unica non rispettata, ovvero la capacità di mantenere l'\vr{Estimate at Completion$^G$} entro la soglia del $\pm 10\%$ rispetto al \vr{Budget at Completion$^G$}.\\
Il picco negativo raggiunto dal gruppo allo Sprint 8 è naturale poiché non era ancora stata iniziata la codifica dell'MVP, e tutte le metriche associate risultavano non accettabili.

\subsection{MQD\_01: Requisiti obbligatori soddisfatti, MQD\_02: Requisiti desiderabili soddisfatti e MQD\_03: Requisiti opzionali soddisfatti}
\begin{figure}[H]
	\centering 
	\includegraphics[width=0.8\textwidth]{PianoQualifica/PB/Requisiti}
	\caption{MQD\_01: Requisiti obbligatori soddisfatti, MQD\_02: Requisiti desiderabili soddisfatti e MQD\_03: Requisiti opzionali soddisfatti}
\end{figure}
Come si nota dal grafico, il gruppo non ha completamente tralasciato i requisiti desiderabili impegnandosi solo sugli obbligatori, ma si è impegnato anche nel raggiungimento di alcuni di essi.\\
La crescita è costante nel tempo e si è raggiunto la totalità dei requisiti obbligatori soddisfatti allo sprint 12, diventando ufficiale nello sprint 13 all'incontro per l'approvazione MVP da parte dell'azienda.

\subsection{MQD\_04: Code Coverage, MQD\_05: Branch Coverage e MQD\_06: Statement Coverage}
\begin{figure}[H]
	\centering 
	\includegraphics[width=0.8\textwidth]{PianoQualifica/PB/Coverage}
	\caption{MQD\_04: Code Coverage, MQD\_05: Branch Coverage e MQD\_06: Statement Coverage}
\end{figure}
Il gruppo è sempre riuscito a tenere un'ottima coverage del progetto.\\
Questo è facilmente dimostrabile considerata la quantità dei test di unità descritta nella sezione \vr{Test di Unità}.

\subsection{MQD\_07: Code Smells per KLOC}
\begin{figure}[H]
	\centering 
	\includegraphics[width=0.8\textwidth]{PianoQualifica/PB/Code_smells}
	\caption{MQD\_07: Code Smells per KLOC}
\end{figure}
Il gruppo ha *riscontrato* dei problemi con questa metrica allo Sprint 9, il primo sprint in cui è iniziata la codifica dell'MVP.\\
Si è subito attivato per risolvere i problemi, sfruttando in maniera più rigorosa le segnalazioni di SonarQube, strumento utilizzato dal gruppo per l'analisi statica.

\subsection{MQD\_08: Cyclomatic Complexity e MQD\_09: Cognitive Complexity}
\begin{figure}[H]
	\centering 
	\includegraphics[width=0.8\textwidth]{PianoQualifica/PB/Complexity}
	\caption{MQD\_08: Cyclomatic Complexity e MQD\_09: Cognitive Complexity}
\end{figure}
Il gruppo ha sempre tenuto sotto controllo la complessità cognitiva e ciclomatica delle funzioni, intervenendo nel caso di funzioni eccessivamente complesse. \\
Grazie alle segnalazioni tempestive di SonarQube, era possibile individuare con precisione dove effettuare il refactoring. \\
Questo approccio si è rivelato molto utile per preservare la manutenibilità del codice nel tempo, anche in ottica di future aggiunte di funzionalità.